\chapter{Introduction}\label{ch:intro}

\section{Problem Background}

The global transition away from fossil fuels requires scalable, carbon-free energy
carriers. Molecular hydrogen~(H\textsubscript{2}) is widely regarded as one of the
most attractive options: it has the highest gravimetric energy density of any known
fuel (120\,MJ/kg), produces only water upon combustion, and can be stored and
transported using existing infrastructure~\cite{turner2004sustainable}.
Currently, however, more than 95\% of global hydrogen is produced by steam methane
reforming --- a process that emits approximately 10\,kg of CO\textsubscript{2} per
kilogram of H\textsubscript{2}~\cite{crabtree2004hydrogen}.

Photoelectrochemical (PEC) water splitting offers a pathway to produce green hydrogen
directly from sunlight and water using a semiconductor photoelectrode immersed in an
aqueous electrolyte~\cite{fujishima1972electrochemical}.
When photons with energy exceeding the semiconductor bandgap are absorbed, they
generate electron--hole pairs.
The holes oxidise water at the photoanode (oxygen evolution reaction, OER), while
the electrons reduce protons at the cathode (hydrogen evolution reaction, HER):

\begin{align}
  \text{Anode (OER):}   &\quad 2\,\mathrm{H_2O}
    \;\xrightarrow{h\nu}\;
    \mathrm{O_2} + 4\,\mathrm{H^+} + 4e^- \\
  \text{Cathode (HER):} &\quad 4\,\mathrm{H^+} + 4e^-
    \;\to\; 2\,\mathrm{H_2}
\end{align}

The overall efficiency of this process is quantified by the Solar-to-Hydrogen (STH)
efficiency:
\begin{equation}
  \mathrm{STH} =
    \frac{J_{ph} \times 1.23\,\mathrm{V} \times \eta_F}{P_{in}} \times 100\%
  \label{eq:sth}
\end{equation}
where $J_{ph}$ is the photocurrent density (mA/cm\textsuperscript{2}),
1.23\,V is the thermodynamic minimum potential for water splitting,
$\eta_F$ is the Faradaic efficiency, and
$P_{in}$ is the incident light power density (mW/cm\textsuperscript{2}).

After decades of research, only a handful of materials come close to the
theoretical STH limit for their bandgap~\cite{walter2010solar}.
Finding a high-efficiency photoelectrode is not a matter of testing one or two
variables --- a researcher working with BiVO\textsubscript{4}, for example,
must decide on the nanostructure morphology, whether to add a co-catalyst,
which electrolyte to use, what pH and applied bias to set, and each choice
interacts with the others in ways that are hard to predict in advance.
Testing all combinations experimentally is simply not feasible within a typical
research timeline.

\section{Motivation and Relevance}

Machine learning has already changed how researchers approach materials discovery
in adjacent fields.
Battery electrolyte screening, catalysis, and photovoltaics all have ML-assisted
workflows now.
PEC water splitting does not --- and the gap is not because the problem is
unimportant, but because the data is hard to collect.

When we went through 529 published papers, we found that fewer than 1 in 10
reported an STH value.
Most papers measure photocurrent and stop there, because calculating STH requires
additional instrumentation and a more careful experimental setup.
The end result is a dataset where the single most important performance metric
practically speaking is also the least frequently occurring one in the training data.
Standard k-fold cross-validation fails quickly when there are only 48 rows.

This project has a genuinely multidisciplinary nature, which ultimately shaped
how I approached the methodology.
Electrochemistry allows you to determine which variables to include as features.
Materials science informs what we can trust in the experimental data.
Machine learning provides the tools to generalise across hundreds of diverse
experiments from different laboratories.
Without all three working together, the result is either a physically invalid model,
or a scientifically valid one that cannot generalise beyond its training data.

\section{Aims and Objectives}

The main goal of this project is to create a practical application of machine
learning.
This involves building a machine-learning based system which predicts three
different measures --- photocurrent density, Solar-to-Hydrogen (STH) efficiency,
and Hydrogen Evolution Rate --- based upon the provided material and experimental
parameters.
Once the model was created, we made the developed system available via a user
interface, allowing researchers to utilise the machine learning based system.

To achieve our stated goals, we needed to solve multiple issues sequentially:

\begin{enumerate}
  \item Create a dataset from 529 peer-reviewed papers representing
        Photoelectrochemical (PEC) water splitting, covering all possible
        combinations of various materials, structural parameters, and operating
        conditions, with photocurrent density, Solar-to-Hydrogen (STH) efficiency,
        and Hydrogen Evolution Rate as predictive variables.
        Unfortunately, creating the above mentioned dataset proved to be far more
        challenging than initially anticipated, because most research studies
        reported only photocurrent density.

  \item Develop numerical representations of input values into features by
        applying established electrochemical relationships.
        These include the bandgap--pH product as described by the Nernst equation,
        an overpotential proxy as derived from Butler--Volmer kinetic equations,
        and hydrogen production rate as derived from Faraday's law.
        In doing so, we avoided requiring the machine learning algorithm to derive
        these same relationships independently.

  \item Train separate regression models utilising XGBoost and Random Forest
        for each predictive variable, evaluate both using 5-fold cross-validation,
        and combine the results into a single weighted ensemble where weights
        represent confidence in each model's output.

  \item Make trained models accessible via a Flask REST API and expose them
        via a React web-based interface, returning predicted values along with
        a LOW / MEDIUM / HIGH classification of potential success based upon
        the input configuration entered by users.

  \item Compare cross-validated performance metrics to demonstrate that adding
        physics-derived estimates and synthetic rows to the existing dataset
        produces significantly improved performance compared to training solely
        on the original sparse data.
\end{enumerate}

\section{Structure of the Report}

The structure of the report is as follows:

\begin{itemize}
  \item In Chapter~\ref{ch:A}, we present all the necessary context to understand
        our design choices: the physical behaviour of PEC cells; which performance
        indicators are important and why; and why a tree-based ensemble model has
        been preferred over a neural network or linear regression.

  \item In Chapter~\ref{ch:B}, the content is focused around the methodical
        development of this project: how the dataset was built from 529 scientific
        articles; how missing values for STH and H\textsubscript{2} were computed
        using physics equations; how 20 features were extracted from 8 input
        parameters; and how the training and prediction pipeline is organised.

  \item Chapter~\ref{ch:C} shows, with examples of three typical material
        configurations, what the models were able to achieve: cross-validation
        results before and after augmentation, feature importance analysis,
        and a system demonstration on representative inputs.

  \item Chapter~\ref{ch:concl} summarises the main outcomes, acknowledges
        where the approach falls short, and outlines the most promising
        directions for future work.
\end{itemize}
